% from ega4-16-fr.tex
\oldpage[IV-4]{11}
\noindent(ii) et (iii) Si $f$ est une immersion, on peut se borner au cas où
$f$ est une immersion fermée, $Y$ étant donc défini par l'Idéal
quasi-cohérent $\mathscr{I}$ de $\mathcal{O}_X$, et $Y^{(n)}$ par l'Idéal
$\mathscr{I}^{n+1}$; les assertions résultent alors de (I, 4.4.5).

Supposons en second lieu que $v$ soit plat; on peut se borner au cas où
$X=\operatorname{Spec}(A)$, $Y=\operatorname{Spec}(B)$,
$X'=\operatorname{Spec}(A')$ sont affines, $A'$ étant un $A$-module plat;
alors $Y'=\operatorname{Spec}(B')$ avec $B'=B\otimes_A A'$; en outre, si
$\mathfrak{I}$ est le noyau de l'homomorphisme $A\to B$, le noyau
$\mathfrak{I}'$ de $A'\to B'$ s'identifie à $\mathfrak{I}\otimes_A A'$ par
platitude, et
\[
  \mathfrak{I}'^n/\mathfrak{I}'^{n+1}
  =(\mathfrak{I}^n/\mathfrak{I}^{n+1})\otimes_A A'.
\]
On en déduit aussitôt, compte tenu de (0\textsubscript{I}, 4.3.3), que le
$\mathcal{O}_{Y'}$-Module
$\mathscr{I}'^n/\mathscr{I}'^{n+1}$ est égal à
\begin{align*}
&\psi'^*\!\left(\sigma^*
  \left((\mathfrak{I}^n/\mathfrak{I}^{n+1})^{\sim}\right)\right)
  \otimes_{\sigma^*(\mathcal{O}_X)}\mathcal{O}_{X'} \\
={}&
\psi'^*\!\left(\sigma^*
  \left((\mathfrak{I}^n/\mathfrak{I}^{n+1})^{\sim}\right)\right)
  \otimes_{\psi'^*(\sigma^*(\mathcal{O}_X))}
  \psi'^*(\mathcal{O}_{X'}) \\
={}&
\rho^*(\mathscr{I}^n/\mathscr{I}^{n+1})
  \otimes_{\rho^*(\psi^*(\mathcal{O}_X))}
  \psi'^*(\mathcal{O}_{X'}),
\end{align*}
et comme en particulier pour $n=0$, on a
\[
  \mathcal{O}_{Y'}=
  \rho^*(\mathcal{O}_Y)
  \otimes_{\rho^*(\psi^*(\mathcal{O}_X))}
  \psi'^*(\mathcal{O}_{X'}),
\]
on obtient un isomorphisme canonique de
$\mathscr{I}'^n/\mathscr{I}'^{n+1}$ sur
\[
  \rho^*(\mathscr{I}^n/\mathscr{I}^{n+1})
  \otimes_{\rho^*(\mathcal{O}_Y)}\mathcal{O}_{Y'}
  =
  (\mathscr{I}^n/\mathscr{I}^{n+1})
  \otimes_{\mathcal{O}_Y}\mathcal{O}_{Y'},
\]
ce qui prouve (iii). Posons maintenant
$C_n=\Gamma(Y,\mathcal{O}_{Y^{(n)}})$,
$C'_n=\Gamma(Y',\mathcal{O}_{Y'^{(n)}})$. Puisque $Y^{(n)}$ et
$Y'^{(n)}$ sont des schémas affines (16.1.5), le noyau
$\mathfrak{K}_n$ (resp. $\mathfrak{K}'_n$) de l'homomorphisme
$C_n\to C_{n-1}$ (resp. $C'_n\to C'_{n-1}$) est
$\Gamma(Y,\mathscr{I}^n/\mathscr{I}^{n+1})$ (resp.
$\Gamma(Y',\mathscr{I}'^n/\mathscr{I}'^{n+1})$); on déduit donc de ce qui
précède que $\mathfrak{K}'_n=\mathfrak{K}_n\otimes_A A'$. Or, on a un
diagramme commutatif
\[
  \xymatrix{
    0 \ar[r] &
    \mathfrak{K}_n\otimes_A A' \ar[r]\ar[d]^r &
    C_n\otimes_A A' \ar[r]\ar[d]^{s_n} &
    C_{n-1}\otimes_A A' \ar[r]\ar[d]^{s_{n-1}} & 0 \\
    0 \ar[r] &
    \mathfrak{K}'_n \ar[r] &
    C'_n \ar[r] &
    C'_{n-1} \ar[r] & 0
  }
\]
où la flèche verticale de gauche est bijective et les deux lignes sont exactes
($A'$ étant un $A$-module plat). On en déduit par récurrence que $s_n$ est
bijectif pour tout $n$, car c'est vrai par hypothèse pour $n=0$, et se déduit
par application du lemme des cinq pour $n$ quelconque. Cela prouve la seconde
assertion de (ii).

\begin{corollary}[16.2.3]
\label{IV.16.2.3}
Soient $g:X\to Y$, $u:Y'\to Y$ deux morphismes de préschémas,
$X'=X\times_Y Y'$, $g':X'\to Y'$ et $v:X'\to X$ les projections canoniques.
Soient $f:Y\to X$ une $Y$-section de $X$ (donc une immersion),
$f'=f_{(Y')}:Y'\to X'$ la $Y'$-section de $X'$ déduite de $f$ par le
changement de base $u$. Alors:
\begin{enumerate}
\item[(i)] Le morphisme
$w_n:Y_{f'}'^{(n)}\to Y_f^{(n)}$ correspondant à $f$, $f'$, $u$, $v$
(16.2.1) et le morphisme canonique
$h'_n:Y_{f'}'^{(n)}\to X'$ identifient $Y_{f'}'^{(n)}$ au produit
$Y_f^{(n)}\times_X X'$.
\item[(ii)] Si l'on munit $\mathcal{O}_{Y_f^{(n)}}$ (resp.
$\mathcal{O}_{Y_{f'}'^{(n)}}$) de la structure de $\mathcal{O}_Y$-Algèbre
définie par $g$ (resp. de la structure de $\mathcal{O}_{Y'}$-Algèbre définie
par $g'$) (16.1.6), l'homomorphisme de $\mathcal{O}_{Y'}$-Algèbres
\[
  \rho^*(\mathcal{O}_{Y_f^{(n)}})
  \otimes_{\mathcal{O}_Y}\mathcal{O}_{Y'}
  \longrightarrow
  \mathcal{O}_{Y_{f'}'^{(n)}}
  \tag{16.2.3.1}
  \label{IV.16.2.3.1}
\]
